% Zumdahl Chapter 16 -- Solubility Equilibria. Elaborated notes, 4 blocks.
% Only SS16.1 is examinable: it holds CED 7.11, 7.12 and 8.11.
\documentclass{apchem-notes}

\apcunitword{Chapter}
\apcunit{16}
\apcunittitle{Solubility Equilibria}
\apcdoctitle{Guided Notes}
\apcsubtitle{Zumdahl \S16.1 \textbullet\ PDF pp.~805--813 \textbullet\ 4 blocks}

\begin{document}
\apctitleblock

\begin{apcnote}[How this chapter fits the AP course]
\textbf{This short chapter completes the CED.} \S16.1 alone carries the last
three topics not covered anywhere else in these materials:
\textbf{7.11} (solubility equilibria and $K_{sp}$), \textbf{7.12} (the
common-ion effect applied to solubility) and \textbf{8.11} (pH and
solubility).

\textbf{Skip the rest of the chapter.} \S16.2 (precipitation and qualitative
analysis) and \S16.3 (complex ion equilibria) are not on the framework.
A short treatment of $Q$ versus $K_{sp}$ closes Block 4 because the
reasoning is Unit 7 material you already own --- but it is flagged as
enrichment, and the chapter test does not assess it.

\textbf{One CED boundary matters more than usual here.} Topic 8.11 carries
an explicit \emph{exclusion statement}: \emph{computations of solubility as
a function of pH will not be assessed}. So pH and solubility is taught
\textbf{qualitatively}, by Le Ch\^atelier, and every quantitative
common-ion problem in these materials uses an ion that is \emph{not}
pH-sensitive. That is a deliberate choice, not an omission.

Nothing here is a new kind of equilibrium. A saturated solution is a
system at equilibrium with a solid, and $K_{sp}$ is its equilibrium
constant with the solid left out --- exactly the rule from Chapter 13.
\end{apcnote}

% =====================================================================
\blocklesson{$K_{sp}$ and the Solubility Product}{Zumdahl \S16.1}

\begin{objectives}
  \item Write the $K_{sp}$ expression for any sparingly soluble salt.
  \item Calculate $K_{sp}$ from a measured solubility.
\end{objectives}

\reading{Zumdahl \S16.1}{805--810}

\begin{warmup}[4]
  \item Species omitted from any $K$ expression:
        \answerblank[6cm]{pure solids and pure liquids}
  \item $K$ for \ce{2A <=> B} in terms of concentrations:
        \answerblank[3.4cm]{$[\ce{B}]/[\ce{A}]^2$}
  \item Is \ce{AgCl} soluble or insoluble by the Unit 4 rules?
        \answerblank[2.4cm]{insoluble}
  \item If $Q < K$, the reaction shifts: \answerblank[2cm]{right}
\end{warmup}

\segment{INSTRUCTION A}{25}{A saturated solution is an equilibrium}

\notesectionz{Dissolving, run backwards}{16.1}
\practice{6}

Put more solid salt into water than can dissolve and two processes run at
once: solid dissolving into ions, and ions returning to the solid. When
their rates are equal the solution is \term{saturated} --- and that is a
\answerblank[4.4cm]{dynamic equilibrium}, not a solution that has ``stopped
dissolving.''

\[ \ce{CaF2(s) <=> Ca^2+(aq) + 2F^-(aq)} \]

The equilibrium constant for this process is the \term{solubility product},
$K_{sp}$:
\[ K_{sp} = [\ce{Ca^2+}][\ce{F^-}]^2 \]

Two features to notice, both of which are old rules reappearing:

\begin{itemize}[leftmargin=*,itemsep=0.3em]
  \item The solid does \emph{not} appear --- it is a
        \answerblank[2.6cm]{pure solid}, exactly as in Chapter 13.
  \item Each ion is raised to the power of its
        \answerblank[2.6cm]{coefficient}, so the 2 in \ce{2F^-} becomes an
        exponent.
\end{itemize}

\begin{apcnote}[Why more solid does not mean more dissolved]
It seems as though a bigger pile of solid --- more surface area --- ought to
give a more concentrated solution. Zumdahl's answer is worth keeping:
doubling the surface area doubles the rate of dissolving, but it
\emph{equally} doubles the rate at which ions re-deposit on that surface.
The two effects cancel and the equilibrium position does not move.

Grinding the solid or stirring the mixture makes equilibrium arrive
\emph{sooner}. Neither changes how much ends up dissolved.
\end{apcnote}

\begin{apctrap}
\textbf{$K_{sp}$ is an equilibrium constant; solubility is an equilibrium
position.} This one sentence organizes the whole chapter.

$K_{sp}$ has a single value for a given solid at a given temperature.
Solubility is how much actually dissolves, and it changes with what else is
in the water --- add a common ion and the solubility falls, while $K_{sp}$
does not move at all. Students who treat the two words as synonyms get
every common-ion question wrong.
\end{apctrap}

\segment{GUIDED PRACTICE}{15}{Writing the expressions}

\begin{enumerate}[leftmargin=*,itemsep=0.35em]
  \item \ce{AgCl}: $K_{sp} = $
        \answerblank[5cm]{$[\ce{Ag+}][\ce{Cl-}]$}
  \item \ce{PbI2}: $K_{sp} = $
        \answerblank[5cm]{$[\ce{Pb^2+}][\ce{I-}]^2$}
  \item \ce{Ag2CrO4}: $K_{sp} = $
        \answerblank[5cm]{$[\ce{Ag+}]^2[\ce{CrO4^2-}]$}
  \item \ce{Mg(OH)2}: $K_{sp} = $
        \answerblank[5cm]{$[\ce{Mg^2+}][\ce{OH-}]^2$}
  \item \ce{Ca3(PO4)2}: $K_{sp} = $
        \answerblank[5cm]{$[\ce{Ca^2+}]^3[\ce{PO4^3-}]^2$}
\end{enumerate}

\segment{INSTRUCTION B}{20}{From measured solubility to $K_{sp}$}

\notesectionz{The ICE table is trivial --- the stoichiometry is not}{16.1}
\practice{5}

If $s$ moles of salt dissolve per litre, the ion concentrations follow
directly from the coefficients. That is the only step where mistakes
happen.

\begin{center}
\begin{tabular}{@{}llll@{}}
\hline
\textbf{Type} & \textbf{Example} & \textbf{Ions at equilibrium} &
  \textbf{$K_{sp}$} \\
\hline
\ce{MX}   & \ce{AgCl}    & $s$, $s$   & $s^2$ \\
\ce{MX2}  & \ce{CaF2}    & $s$, $2s$  & $(s)(2s)^2 = 4s^3$ \\
\ce{M2X}  & \ce{Ag2CrO4} & $2s$, $s$  & $(2s)^2(s) = 4s^3$ \\
\ce{M3X}  & \ce{Ag3PO4}  & $3s$, $s$  & $(3s)^3(s) = 27s^4$ \\
\hline
\end{tabular}
\end{center}

Note that \ce{MX2} and \ce{M2X} give the \emph{same} formula, $4s^3$ ---
which is why they can be compared with each other but not with an \ce{MX}
salt.

\begin{workedexample}[Worked example 1: Zumdahl's copper(I) bromide]
\ce{CuBr} has a measured solubility of
$2.0\times10^{-4}$~mol/L at \SI{25}{\celsius}. Find $K_{sp}$.

\[ \ce{CuBr(s) <=> Cu+(aq) + Br^-(aq)} \qquad
   K_{sp} = [\ce{Cu+}][\ce{Br-}] \]
Each formula unit gives one of each ion, so
$[\ce{Cu+}] = [\ce{Br-}] = 2.0\times10^{-4}$:
\[ K_{sp} = (2.0\times10^{-4})(2.0\times10^{-4})
   = \mathbf{4.0\times10^{-8}} \]
Units are conventionally omitted.
\end{workedexample}

\begin{workedexample}[Worked example 2: a 1:2 salt, and grams to moles]
Lead(II) iodide dissolves to the extent of \SI{0.70}{\gram\per\litre} at
\SI{25}{\celsius}. Find $K_{sp}$. ($M = \SI{461.0}{\gram\per\mole}$)

\textbf{Convert first.} $s = 0.70/461.0 = 1.52\times10^{-3}$~mol/L.

\textbf{Then the stoichiometry.}
$\ce{PbI2(s) <=> Pb^2+(aq) + 2I^-(aq)}$, so $[\ce{Pb^2+}] = s$ and
$[\ce{I-}] = 2s = 3.04\times10^{-3}$:
\[ K_{sp} = (1.52\times10^{-3})(3.04\times10^{-3})^2
   = \mathbf{1.4\times10^{-8}} \]

Forgetting to double the iodide gives $3.5\times10^{-9}$ --- low by a
factor of 4, and a guaranteed lost point.
\end{workedexample}

\segment{APPLICATION}{20}{Calculating $K_{sp}$}

\begin{enumerate}[leftmargin=*,itemsep=0.4em]
  \item \ce{BaSO4} dissolves to $9.0\times10^{-3}$~g/L
        ($M = \SI{233.4}{\gram\per\mole}$). Find $K_{sp}$.
        \workspace[2.6cm]
        \keyonly{\begin{apcanswer}$s = 9.0\times10^{-3}/233.4 =
        3.86\times10^{-5}$~M; both ions equal $s$;
        $K_{sp} = (3.86\times10^{-5})^2 =
        \mathbf{1.5\times10^{-9}}$\end{apcanswer}}
  \item \ce{Ag2CrO4} has a solubility of $1.3\times10^{-4}$~M. Find
        $K_{sp}$. \workspace[2.6cm]
        \keyonly{\begin{apcanswer}$[\ce{Ag+}] = 2s = 2.6\times10^{-4}$,
        $[\ce{CrO4^2-}] = s = 1.3\times10^{-4}$;
        $K_{sp} = (2.6\times10^{-4})^2(1.3\times10^{-4}) =
        \mathbf{8.8\times10^{-12}}$\end{apcanswer}}
  \item A student computes $K_{sp}$ for \ce{CaF2} as $(s)(s)^2$. What did
        they miss, and by what factor is the answer wrong?
        \answerlines[2]{They failed to double the fluoride:
        $[\ce{F-}] = 2s$, not $s$. The correct expression is
        $(s)(2s)^2 = 4s^3$, so their answer is too small by a factor
        of 4.}
\end{enumerate}

\begin{exitticket}
A saturated solution of \ce{AgCl} sits over a pile of undissolved
\ce{AgCl}. A student adds a second spoonful of solid \ce{AgCl}. What
happens to $[\ce{Ag+}]$, and what happens to $K_{sp}$?
\keyonly{\begin{apcanswer}Neither changes. The solid does not appear in the
$K_{sp}$ expression, and adding more of it changes the rates of dissolving
and re-depositing equally, so the equilibrium position does not move.
$K_{sp}$ depends only on temperature.\end{apcanswer}}
\end{exitticket}

% =====================================================================
\blocklesson{Solubility from $K_{sp}$, and Comparing Salts}{Zumdahl \S16.1}

\begin{objectives}
  \item Calculate molar solubility from $K_{sp}$ for any stoichiometry.
  \item Judge when $K_{sp}$ values may be compared directly and when they
        may not.
\end{objectives}

\reading{Zumdahl \S16.1 (Relative Solubilities)}{808--811}

\begin{warmup}[3]
  \item $K_{sp}$ for an \ce{MX} salt in terms of $s$:
        \answerblank[2cm]{$s^2$}
  \item $K_{sp}$ for an \ce{MX2} salt in terms of $s$:
        \answerblank[2cm]{$4s^3$}
  \item $K_{sp}$ is a constant; solubility is a
        \answerblank[7.6cm]{position (it varies with the solution)}
\end{warmup}

\segment{INSTRUCTION A}{25}{Running the calculation backwards}

\notesectionz{Solving for $s$}{16.1}
\practice{5}

Going from $K_{sp}$ to solubility is the same relationship read the other
way. Rearranged:

\begin{center}
\begin{tabular}{@{}lll@{}}
\hline
\textbf{Type} & \textbf{Relation} & \textbf{Solubility} \\
\hline
\ce{MX}  & $K_{sp} = s^2$   & $s = \sqrt{K_{sp}}$ \\
\ce{MX2} or \ce{M2X} & $K_{sp} = 4s^3$ &
  $s = \sqrt[3]{K_{sp}/4}$ \\
\ce{M3X} & $K_{sp} = 27s^4$ & $s = \sqrt[4]{K_{sp}/27}$ \\
\hline
\end{tabular}
\end{center}

\begin{workedexample}[Worked example 3: two stoichiometries side by side]
\textbf{\ce{AgCl}}, $K_{sp} = 1.6\times10^{-10}$.
Let $s = [\ce{Ag+}] = [\ce{Cl-}]$:
\[ 1.6\times10^{-10} = s^2 \;\Longrightarrow\;
   s = \mathbf{1.3\times10^{-5}}~\text{M} \]

\textbf{\ce{CaF2}}, $K_{sp} = 4.0\times10^{-11}$.
Here $[\ce{Ca^2+}] = s$ and $[\ce{F-}] = 2s$:
\[ 4.0\times10^{-11} = (s)(2s)^2 = 4s^3
   \;\Longrightarrow\; s^3 = 1.0\times10^{-11}
   \;\Longrightarrow\; s = \mathbf{2.2\times10^{-4}}~\text{M} \]

Notice: \ce{CaF2} has the \emph{smaller} $K_{sp}$ but is over ten times
\emph{more} soluble. Hold that thought.
\end{workedexample}

\segment{INSTRUCTION B}{20}{When you may compare $K_{sp}$ values}

\notesectionz{The comparison trap}{16.1}
\practice{6}

\begin{center}
\fbox{\parbox{0.88\linewidth}{\centering
$K_{sp}$ values may be compared directly \textbf{only} for salts that
produce the \textbf{same total number of ions}.}}
\end{center}

\textbf{Same stoichiometry --- comparison works.} All three of these are
\ce{MX} salts, so $s = \sqrt{K_{sp}}$ for each and the order of $K_{sp}$ is
the order of solubility:

\begin{center}
\begin{tabular}{@{}lcc@{}}
\hline
\textbf{Salt} & $K_{sp}$ & \textbf{solubility (M)} \\
\hline
\ce{AgI}   & $1.5\times10^{-16}$ & $1.2\times10^{-8}$ \\
\ce{CuI}   & $5.0\times10^{-12}$ & $2.2\times10^{-6}$ \\
\ce{CaSO4} & $6.1\times10^{-5}$  & $7.8\times10^{-3}$ \\
\hline
\end{tabular}
\end{center}

\textbf{Different stoichiometry --- the order can reverse outright.}
Zumdahl's example is worth memorizing because it is so extreme:

\begin{center}
\begin{tabular}{@{}lccc@{}}
\hline
\textbf{Salt} & \textbf{type} & $K_{sp}$ & \textbf{solubility (M)} \\
\hline
\ce{CuS}   & \ce{MX}    & $8.5\times10^{-45}$ & $9.2\times10^{-23}$ \\
\ce{Ag2S}  & \ce{M2X}   & $1.6\times10^{-49}$ & $3.4\times10^{-17}$ \\
\ce{Bi2S3} & \ce{M2X3}  & $1.1\times10^{-73}$ & $1.0\times10^{-15}$ \\
\hline
\end{tabular}
\end{center}

$K_{sp}$ order: \ce{CuS} $>$ \ce{Ag2S} $>$ \ce{Bi2S3}. \\
Solubility order: \answerblank[6cm]{\ce{Bi2S3} $>$ \ce{Ag2S} $>$ \ce{CuS}}
--- \textbf{exactly reversed}.

\begin{apctrap}
The reversal is not mysterious. \ce{Bi2S3} produces \emph{five} ions per
formula unit, so its $K_{sp}$ is a product of five concentration factors,
each tiny. Multiplying five small numbers gives an absurdly small
$K_{sp}$ even when a respectable amount of solid has dissolved.

The rule to carry into the exam: \textbf{if the formulas do not have the
same ion count, you must calculate $s$ before you may rank them.}
\end{apctrap}

\notesectionz{Connecting back to the Unit 4 solubility rules}{16.1}

The solubility rules you memorized in Unit 4 are the qualitative version of
this same idea. The CED makes the link explicit: a salt with
$K_{sp} > 1$ is what those rules call
\answerblank[2.4cm]{soluble}, and the very small $K_{sp}$ values in the
table belong to the salts the rules call
\answerblank[2.4cm]{insoluble}. ``Insoluble'' never meant zero --- it meant
an equilibrium lying far to the left.

\segment{APPLICATION}{20}{Ranking and calculating}

\begin{enumerate}[leftmargin=*,itemsep=0.4em]
  \item Calculate the molar solubility of \ce{Mg(OH)2},
        $K_{sp} = 8.9\times10^{-12}$. \workspace[2.4cm]
        \keyonly{\begin{apcanswer}$[\ce{Mg^2+}] = s$, $[\ce{OH-}] = 2s$;
        $8.9\times10^{-12} = 4s^3$; $s^3 = 2.2\times10^{-12}$;
        $s = \mathbf{1.3\times10^{-4}}$~M\end{apcanswer}}
  \item Rank \ce{BaSO4} ($1.5\times10^{-9}$), \ce{SrSO4}
        ($3.2\times10^{-7}$) and \ce{PbSO4} ($1.3\times10^{-8}$) by
        increasing solubility, and say why no calculation is needed.
        \answerlines[2]{\ce{BaSO4} $<$ \ce{PbSO4} $<$ \ce{SrSO4}. All three
        are \ce{MX} salts producing two ions, so $s = \sqrt{K_{sp}}$ and the
        $K_{sp}$ order is already the solubility order.}
  \item \ce{CaF2} ($K_{sp} = 4.0\times10^{-11}$) and \ce{BaCrO4}
        ($K_{sp} = 8.5\times10^{-11}$). Which is more soluble?
        \workspace[2.6cm]
        \keyonly{\begin{apcanswer}Calculation is required --- different ion
        counts. \ce{CaF2}: $s = \sqrt[3]{4.0\times10^{-11}/4} =
        2.2\times10^{-4}$. \ce{BaCrO4}: $s = \sqrt{8.5\times10^{-11}} =
        9.2\times10^{-6}$. \textbf{\ce{CaF2} is about 24 times more
        soluble}, despite the smaller $K_{sp}$.\end{apcanswer}}
\end{enumerate}

\begin{exitticket}
Salt A has a larger $K_{sp}$ than salt B. A student concludes A is more
soluble. Under what condition are they right, and under what condition
could they be badly wrong?
\keyonly{\begin{apcanswer}They are right if A and B produce the same total
number of ions, since then solubility is the same increasing function of
$K_{sp}$ for both. They can be wrong whenever the ion counts differ --- a
salt producing more ions can have a far smaller $K_{sp}$ and still be much
more soluble, as \ce{Bi2S3} is compared with \ce{CuS}.\end{apcanswer}}
\end{exitticket}

% =====================================================================
\blocklesson{The Common-Ion Effect on Solubility}{Zumdahl \S16.1}

\begin{objectives}
  \item Predict qualitatively how a common ion changes solubility.
  \item Calculate solubility in a solution already containing a common ion.
\end{objectives}

\reading{Zumdahl \S16.1 (Common Ion Effect)}{811--812}

\begin{warmup}[4]
  \item Molar solubility of \ce{AgCl} in pure water
        ($K_{sp} = 1.6\times10^{-10}$):
        \answerblank[3cm]{$1.3\times10^{-5}$~M}
  \item Adding a product shifts an equilibrium:
        \answerblank[2cm]{left}
  \item A common ion is one that is:
        \answerblank[5.4cm]{already in the equilibrium}
  \item Does adding a common ion change $K_{sp}$?
        \answerblank[2cm]{no}
\end{warmup}

\segment{INSTRUCTION A}{25}{Le Ch\^atelier, one more time}

\notesectionz{Why solubility falls}{16.1}
\practice{6}

\[ \ce{CaF2(s) <=> Ca^2+(aq) + 2F^-(aq)} \]

Dissolve the salt in a solution that already contains \ce{F-} --- from
\ce{NaF}, say --- and you have added a \answerblank[2cm]{product}. The
equilibrium shifts \answerblank[2cm]{left}, so \emph{less} \ce{CaF2}
dissolves.

This is the identical mechanism you met in Chapter 15, where \ce{NaF}
suppressed the ionization of \ce{HF}. Only the equilibrium has changed.

\begin{apctrap}
Say precisely what moved. The \emph{solubility} decreased; $K_{sp}$ did
\textbf{not} change --- it cannot, because nothing about the temperature
changed. The product $[\ce{Ca^2+}][\ce{F-}]^2$ still equals
$4.0\times10^{-11}$; the two concentrations have simply redistributed, with
much more fluoride and much less calcium.

``The common ion lowers $K_{sp}$'' is marked wrong every time.
\end{apctrap}

\begin{workedexample}[Worked example 4: Zumdahl's calcium fluoride]
Find the solubility of \ce{CaF2} ($K_{sp} = 4.0\times10^{-11}$) in
\SI{0.025}{\Molar} \ce{NaF}, and compare with pure water.

Let $s$ be the moles of \ce{CaF2} that dissolve per litre.

\begin{center}
\begin{tabular}{@{}lcc@{}}
\hline
 & $[\ce{Ca^2+}]$ & $[\ce{F-}]$ \\
\hline
initial     & 0 & 0.025 \\
change      & $+s$ & $+2s$ \\
equilibrium & $s$ & $0.025 + 2s$ \\
\hline
\end{tabular}
\end{center}

\[ 4.0\times10^{-11} = (s)(0.025 + 2s)^2 \approx (s)(0.025)^2 \]
because $K_{sp}$ is tiny, so $2s$ is negligible beside 0.025. Then
\[ s = \frac{4.0\times10^{-11}}{6.25\times10^{-4}}
     = \mathbf{6.4\times10^{-8}}~\text{M} \]

In pure water the solubility was $2.2\times10^{-4}$~M. The fluoride already
present cut it by a factor of about \textbf{3400}.
\end{workedexample}

\segment{GUIDED PRACTICE}{15}{Which ion, and which power}

The size of the effect depends on \emph{which} ion you add, because the two
ions enter the expression with different exponents.

\begin{enumerate}[leftmargin=*,itemsep=0.35em]
  \item \ce{AgCl} in \SI{0.10}{\Molar} \ce{NaCl}: $s = K_{sp}/[\ce{Cl-}] =$
        \answerblank[3.2cm]{$1.6\times10^{-9}$~M}
  \item \ce{AgCl} in \SI{0.020}{\Molar} \ce{AgNO3}:
        \answerblank[3.2cm]{$8.0\times10^{-9}$~M}
  \item \ce{BaSO4} ($1.5\times10^{-9}$) in \SI{0.010}{\Molar}
        \ce{Na2SO4}: \answerblank[3.2cm]{$1.5\times10^{-7}$~M}
  \item Would \ce{NaNO3} change the solubility of \ce{AgCl}? Why?
        \answerblank[8cm]{no --- neither ion is in the equilibrium}
\end{enumerate}

\segment{INSTRUCTION B}{20}{Adding the cation instead}

\notesectionz{The exponent decides the sensitivity}{16.1}
\practice{5}

For \ce{CaF2}, adding \ce{F-} is far more effective than adding
\ce{Ca^2+}, because fluoride is \answerblank[2.2cm]{squared} in the
expression while calcium is not.

\begin{workedexample}[Worked example 5: the same salt, the other ion]
Find the solubility of \ce{CaF2} in \SI{0.10}{\Molar} \ce{Ca(NO3)2}.

Now the \emph{calcium} is pre-loaded: $[\ce{Ca^2+}] = 0.10 + s \approx 0.10$
and $[\ce{F-}] = 2s$.
\[ 4.0\times10^{-11} = (0.10)(2s)^2 = 0.40s^2
   \;\Longrightarrow\; s^2 = 1.0\times10^{-10}
   \;\Longrightarrow\; s = \mathbf{1.0\times10^{-5}}~\text{M} \]

Note that the \emph{fluoride is still doubled} even though calcium is the
ion you added --- the $(2s)^2$ contributes a factor of 4. Writing
$(0.10)(s)^2$ here is the commonest slip in the whole chapter.

\textbf{Compare the three cases:}

\begin{center}
\begin{tabular}{@{}lcc@{}}
\hline
\textbf{Solvent} & \textbf{solubility (M)} & \textbf{fall} \\
\hline
pure water                  & $2.2\times10^{-4}$ & --- \\
\SI{0.10}{\Molar} \ce{Ca(NO3)2} & $1.0\times10^{-5}$ & $\times 22$ \\
\SI{0.10}{\Molar} \ce{NaF}      & $4.0\times10^{-9}$ & $\times 54{,}000$ \\
\hline
\end{tabular}
\end{center}

Same concentration of common ion, wildly different effect --- because the
fluoride is squared.
\end{workedexample}

\segment{APPLICATION}{20}{Common-ion calculations}

\begin{enumerate}[leftmargin=*,itemsep=0.4em]
  \item Calculate the solubility of \ce{PbI2}
        ($K_{sp} = 1.4\times10^{-8}$) in \SI{0.10}{\Molar} \ce{NaI}.
        \workspace[2.6cm]
        \keyonly{\begin{apcanswer}$[\ce{I-}] \approx 0.10$;
        $1.4\times10^{-8} = (s)(0.10)^2$;
        $s = \mathbf{1.4\times10^{-6}}$~M --- about 1100 times less than
        in pure water\end{apcanswer}}
  \item Explain, without calculating, why \ce{BaSO4} is used as an X-ray
        contrast agent despite barium ions being toxic.
        \answerlines[3]{$K_{sp}$ is very small ($1.5\times10^{-9}$), so
        almost no \ce{Ba^2+} enters solution. The suspension is also
        administered with excess sulfate present, and by the common-ion
        effect that suppresses the dissolution further still.}
  \item A student says a common ion works by ``using up'' the solid.
        Correct them. \answerlines[2]{Nothing is used up --- the common ion
        shifts the position of the dissolution equilibrium to the left, so
        less solid dissolves in the first place. More undissolved solid
        remains, not less.}
\end{enumerate}

\begin{exitticket}
\ce{CaF2} is dissolved to saturation in \SI{0.10}{\Molar} \ce{NaF}. State
what happens to each of: the solubility of \ce{CaF2}, $[\ce{F-}]$,
$[\ce{Ca^2+}]$, and $K_{sp}$ --- compared with pure water.
\keyonly{\begin{apcanswer}Solubility: much lower. $[\ce{F-}]$: much higher
(it is set by the \ce{NaF}). $[\ce{Ca^2+}]$: much lower, since it equals the
solubility. $K_{sp}$: unchanged --- the product of the two still equals
$4.0\times10^{-11}$.\end{apcanswer}}
\end{exitticket}

% =====================================================================
\blocklesson{pH and Solubility}{Zumdahl \S16.1--16.2}

\begin{objectives}
  \item Predict qualitatively how pH changes the solubility of a salt.
  \item Use $Q$ against $K_{sp}$ to predict precipitation \emph{(enrichment)}.
\end{objectives}

\reading{Zumdahl \S16.1 (pH and Solubility)}{812--813}

\begin{warmup}[3]
  \item Common ion added $\Rightarrow$ solubility:
        \answerblank[2cm]{falls}
  \item Conjugate base of a \emph{weak} acid is a:
        \answerblank[5.6cm]{meaningful (effective) base}
  \item Conjugate base of a \emph{strong} acid is a:
        \answerblank[4.4cm]{negligible base}
\end{warmup}

\begin{apcnote}[Read this before Block 4]
CED topic 8.11 carries an \textbf{exclusion statement}:
\emph{computations of solubility as a function of pH will not be assessed on
the AP Exam.} Everything in this block is therefore \textbf{qualitative} ---
predict the direction and justify it with Le Ch\^atelier. If a question ever
asks you to \emph{calculate} a solubility at a given pH, it is not an AP
question.
\end{apcnote}

\segment{INSTRUCTION A}{25}{When does pH matter?}

\notesectionz{Only if the anion is a base}{16.1}
\practice{6}

\[ \ce{Mg(OH)2(s) <=> Mg^2+(aq) + 2OH^-(aq)} \]

Add acid. The \ce{H+} reacts with \ce{OH-} and removes it from solution.
That is removing a \answerblank[2cm]{product}, so the equilibrium shifts
\answerblank[2.2cm]{right} and \emph{more} solid dissolves.

This is why milk of magnesia --- a suspension of solid \ce{Mg(OH)2} ---
dissolves in the stomach exactly as fast as it is needed to neutralize
excess acid.

Now the general principle. Adding acid can only pull an anion out of
solution if that anion actually reacts with \ce{H+} --- that is, if it is a
\term{meaningful base}, which happens when its conjugate acid \ce{HX} is
\answerblank[2.4cm]{weak}.

\begin{center}
\begin{tabular}{@{}>{\raggedright\arraybackslash}p{4.2cm}%
                 >{\raggedright\arraybackslash}p{3.8cm}%
                 >{\raggedright\arraybackslash}p{4.8cm}@{}}
\hline
\textbf{Anion} & \textbf{Conjugate acid} & \textbf{Solubility in acid} \\
\hline
\ce{OH-}, \ce{S^2-}, \ce{CO3^2-}, \ce{C2O4^2-}, \ce{CrO4^2-},
  \ce{PO4^3-}, \ce{F-}
  & \emph{weak} --- \ce{H2O}, \ce{H2S}, \ce{H2CO3}, \ce{HF}
  & \textbf{increases}: the anion is consumed \\[0.7em]
\ce{Cl-}, \ce{Br-}, \ce{I-}, \ce{NO3-}, \ce{ClO4-}
  & \emph{strong} --- \ce{HCl}, \ce{HNO3}
  & \textbf{unchanged}: the anion does not react \\
\hline
\end{tabular}
\end{center}

\begin{workedexample}[Worked example 6: three salts, three answers]
\textbf{\ce{Ag3PO4} in acid --- more soluble.} \ce{PO4^3-} is a strong
enough base to take a proton:
\[ \ce{H+(aq) + PO4^3-(aq) -> HPO4^2-(aq)} \]
That lowers $[\ce{PO4^3-}]$, so
$\ce{Ag3PO4(s) <=> 3Ag+(aq) + PO4^3-(aq)}$ shifts right.

\textbf{\ce{AgCl} in acid --- unchanged.} \ce{Cl-} is the conjugate base of
\ce{HCl}, a strong acid, so it has essentially no affinity for protons. No
\ce{HCl} molecules form, $[\ce{Cl-}]$ is untouched, and the equilibrium does
not move.

\textbf{\ce{Mg(OH)2} in \emph{base} --- less soluble.} Raising the pH adds
\ce{OH-}, which is a product. This is the ordinary common-ion effect, and
the solubility falls.
\end{workedexample}

\begin{apctrap}
``Acids are reactive, so everything dissolves better in acid'' is false, and
Zumdahl asks you to refute it directly. \ce{AgCl} is the counterexample:
its solubility in \SI{1}{\Molar} \ce{HNO3} is the same as in pure water.
The test is never how strong the acid is --- it is whether the
\emph{anion} is a base.
\end{apctrap}

\segment{GUIDED PRACTICE}{15}{More soluble in acid?}

For each salt, write \textbf{yes} or \textbf{no} and name the anion's
conjugate acid:

\begin{enumerate}[leftmargin=*,itemsep=0.3em]
  \item \ce{CaCO3}: \answerblank[6cm]{yes --- \ce{H2CO3} is weak}
  \item \ce{AgBr}: \answerblank[6cm]{no --- \ce{HBr} is strong}
  \item \ce{CaF2}: \answerblank[6cm]{yes --- \ce{HF} is weak}
  \item \ce{Fe(OH)3}: \answerblank[6cm]{yes --- \ce{H2O}, very weak}
  \item \ce{PbI2}: \answerblank[6cm]{no --- \ce{HI} is strong}
\end{enumerate}

\notesectionz{Where limestone caves come from}{16.1}

Rainwater dissolves atmospheric \ce{CO2}, which makes it slightly acidic.
That acid attacks limestone:
\[ \ce{CaCO3(s) + H+(aq) -> Ca^2+(aq) + HCO3^-(aq)} \]
Over long times this hollows out caverns. Where the water drips into open
air, dissolved \ce{CO2} escapes, the pH rises, and the reverse happens ---
\ce{CaCO3} precipitates as stalactites and stalagmites. The same
equilibrium, run in both directions by nothing more than a change in pH.

\segment{APPLICATION}{20}{Predicting and justifying}

\begin{enumerate}[leftmargin=*,itemsep=0.4em]
  \item Predict the effect of lowering the pH on the solubility of
        \ce{ZnS}, and justify with an equation.
        \answerlines[3]{Solubility increases. \ce{S^2-} is the conjugate
        base of the weak acid \ce{H2S}, so
        $\ce{S^2-(aq) + H+(aq) -> HS^-(aq)}$ removes sulfide from solution.
        By Le Ch\^atelier, $\ce{ZnS(s) <=> Zn^2+ + S^2-}$ shifts right.}
  \item A solution is saturated with \ce{Mg(OH)2}. \ce{NaOH} is added.
        State the effect on solubility, on $[\ce{Mg^2+}]$, and on
        $K_{sp}$. \answerlines[2]{Solubility falls and $[\ce{Mg^2+}]$
        falls, by the common-ion effect. $K_{sp}$ is unchanged --- only
        temperature changes it.}
  \item Two beakers hold saturated \ce{AgCl} and saturated \ce{CaCO3}.
        Nitric acid is added to both. Describe and explain what differs.
        \answerlines[3]{The \ce{CaCO3} dissolves further, because
        \ce{CO3^2-} is a meaningful base and is consumed by \ce{H+},
        shifting its equilibrium right. The \ce{AgCl} is unaffected,
        because \ce{Cl-} is the conjugate base of a strong acid and does
        not react with \ce{H+}.}
\end{enumerate}

\segment{INSTRUCTION B}{20}{$Q$ versus $K_{sp}$ \emph{(enrichment)}}

\begin{apcnote}[Enrichment --- not on the chapter test]
Predicting precipitation is Zumdahl \S16.2, which is off the CED as a topic.
It is included here because the reasoning is nothing new: it is the
$Q$-versus-$K$ comparison from Unit 7, applied to a dissolution.
\end{apcnote}

Compute the \term{ion product} $Q$ using the concentrations \emph{as mixed},
before any reaction, then compare:

\begin{center}
\begin{tabular}{@{}ll@{}}
\hline
$Q > K_{sp}$ & a precipitate forms, until the ions satisfy $K_{sp}$ \\
$Q = K_{sp}$ & exactly saturated; no change \\
$Q < K_{sp}$ & no precipitate; the solution is unsaturated \\
\hline
\end{tabular}
\end{center}

\begin{workedexample}[Worked example 7: does a precipitate form?]
\SI{50.0}{\milli\litre} of \SI{0.0020}{\Molar} \ce{AgNO3} is mixed with
\SI{50.0}{\milli\litre} of \SI{0.0020}{\Molar} \ce{NaCl}.
$K_{sp}(\ce{AgCl}) = 1.6\times10^{-10}$.

\textbf{Dilution first} --- the total volume is now
\SI{100.0}{\milli\litre}, so each concentration is halved:
$[\ce{Ag+}]_0 = [\ce{Cl-}]_0 = 1.0\times10^{-3}$~M.

\[ Q = (1.0\times10^{-3})(1.0\times10^{-3}) = 1.0\times10^{-6} \]
$Q = 1.0\times10^{-6}$ is far greater than
$K_{sp} = 1.6\times10^{-10}$, so \textbf{a precipitate forms}.

Skipping the dilution step is the usual error --- and it changes $Q$ by a
factor of 4 here.
\end{workedexample}

\begin{exitticket}
Give the single test that decides whether a salt's solubility depends on
pH, and apply it to \ce{BaSO4} and to \ce{BaCO3}.
\keyonly{\begin{apcanswer}Ask whether the anion is a meaningful base ---
equivalently, whether its conjugate acid is weak. \ce{SO4^2-} comes from
\ce{H2SO4}, essentially strong, so \ce{BaSO4} is insensitive to pH.
\ce{CO3^2-} comes from the weak acid \ce{H2CO3}, so \ce{BaCO3} dissolves
more readily as the pH falls.\end{apcanswer}}
\end{exitticket}

\end{document}
